\documentclass[11pt,a4paper]{article}
\usepackage[margin=2.5cm]{geometry}
\usepackage{graphicx}
\usepackage{amsmath,amssymb}
\usepackage{booktabs}
\usepackage{hyperref}
\usepackage{bookmark}
\usepackage{xcolor}
\usepackage{float}
\usepackage{tabularx}
\usepackage{fontspec}
\usepackage{ucharclasses}
\usepackage{multirow}
\usepackage{array}
\usepackage{longtable}

\setmainfont{Noto Sans CJK KR}
\newfontfamily\hangulfont{Noto Sans CJK KR}
\setTransitionsForCJK{\hangulfont}{}{}

\tolerance=1000
\emergencystretch=3em
\hfuzz=2pt

\definecolor{success}{HTML}{2E7D32}
\definecolor{failure}{HTML}{C62828}
\definecolor{partial}{HTML}{F57F17}
\definecolor{inconclusive}{HTML}{757575}
\definecolor{tableheader}{HTML}{E3F2FD}

\newcommand{\teff}{T_\mathrm{eff}}
\newcommand{\Cq}{C(q)}
\newcommand{\knorm}{\|k\|^2}

\title{Boltzmann Thermodynamics of Attention과 FOKVQ:\\[0.3em]
\large 10개 실험 종합 보고서\\[0.3em]
\small{Definitive Edition --- Critic Rounds 1--5 Corrections Applied}}
\author{2026-03-30}
\date{}

\begin{document}
\maketitle

%% ==========================================================================
%% Abstract
%% ==========================================================================
\begin{abstract}
10개 실험(1-1 -- 3-3)을 세 가지 중심 발견으로 재구성하여 보고한다.
첫째, PCA 기반 양자화 축 선택은 LDA(K MSE 26--30배), AFOD($p > 0.4$), Lie Group 수치 최적화(4.3배)를 포함한 모든 대안에 대해 압도적 최적성을 보였다.
둘째, K-space facet 구조는 코퍼스 간 유의하게 분리되나(gap $= 0.408$, $p = 1.6 \times 10^{-64}$), 양자화 축 선택에는 정적 PCA로 충분하다.
셋째, Boltzmann 유효 온도 프레임워크는 표준 MHA(GPT-2)에서 점진적 냉각을 포착하나($\rho = -0.657$), GQA에서 파괴되고 비열 기반 양자화가 PPL을 개선하지 못하여 기술적(descriptive) 가치에 한정된다.
\end{abstract}

\tableofcontents
\newpage

%% ==========================================================================
%% Section 1: Executive Summary
%% ==========================================================================
\section{Executive Summary}

이 절은 10개 실험의 판정(verdict)과 핵심 수치를 단일 테이블로 제시하여, 읽는 이가 전체 실험 시리즈의 결론을 신속히 파악할 수 있도록 한다.

\subsection{종합 판정 매트릭스}

Table~\ref{tab:exec_verdict}는 10개 실험의 핵심 지표, 판정, 함의를 요약한다.
모든 수치는 해당 실험의 결과 JSON에서 직접 추출하였다.

\begin{table}[H]
\centering
\caption{10개 실험 종합 판정 매트릭스.}
\label{tab:exec_verdict}
\small
\begin{tabularx}{\textwidth}{@{}ll>{\raggedright\arraybackslash}p{2.2cm}>{\raggedright\arraybackslash}X@{}}
\toprule
\textbf{Exp} & \textbf{Key Metric} & \textbf{Verdict} & \textbf{Implication} \\
\midrule
1-1a & $\rho = -0.657$, $p = 4.82 \times 10^{-4}$ & \textcolor{partial}{Weak Success} & MHA progressive cooling 확인 \\
1-1b & $\rho = +0.108$, $p = 0.530$ & \textcolor{failure}{Failure} & GQA가 cooling 파괴 \\
1-2 & disagree $= 39.8\%$ & \textcolor{partial}{Not Better} & $\Cq$는 $\knorm$와 독립적이나 PPL 미개선 \\
1-3 & $\Delta H = -0.122$, $p < 10^{-10}$ & \textcolor{failure}{Rejected} & FFN 주파수 선택성 기각 \\
2-1 & gap $= 0.408$, $p = 1.6 \times 10^{-64}$ & \textcolor{success}{Supported} & Facet K-space 코퍼스 간 분리 \\
2-2 & PCA K MSE 26--30$\times$ & \textcolor{success}{PCA Wins} & PCA가 LDA를 압도 \\
2-3 & 구현 A/B 상충 & \textcolor{inconclusive}{Inconclusive} & 재실험 필요 \\
2-4 & AFOD $\approx$ PCA ($p > 0.4$) & \textcolor{success}{PCA Sufficient} & AFOD 초기화 불필요 \\
3-1 & PCA $4.3\times$ opt. & \textcolor{success}{PCA Near-Opt} & SO($d$) 최적화가 PCA 미달 \\
3-2 & N=256 NMSE $= 0.157$ & \textcolor{failure}{Not Viable} & SSM distillation에 큰 state dim 필요 \\
3-3 & RoPE best (NMSE $= 0.475$) & \textcolor{failure}{Spectral Not Best} & Spectral 초기화 가설 기각 \\
\bottomrule
\end{tabularx}
\end{table}

Table~\ref{tab:exec_verdict}에서 확인할 수 있듯이, 양자화 축 선택에 관한 실험(2-2, 2-4, 3-1)은 일관되게 PCA의 최적성을 지지한다.
Boltzmann 관련 실험(1-1 -- 1-3)은 현상 기술에는 성공하였으나 규범적 활용에는 실패하였고, SSM 관련 실험(3-2, 3-3)은 attention 패턴의 SSM 압축이 현 조건에서 비실용적임을 보여준다.

%% ==========================================================================
%% Section 2: Central Finding 1 --- PCA의 보편적 최적성
%% ==========================================================================
\section{Central Finding 1: PCA의 보편적 최적성}

이 절은 양자화 축 선택에서 PCA가 모든 대안을 압도한다는 본 실험 시리즈의 가장 강력한 발견을 종합한다.
네 개의 독립적 실험(2-2, 2-4, 3-1, 3-2)이 서로 다른 각도에서 동일한 결론을 지지한다.

\subsection{PCA vs LDA: K MSE 26--30배 우위 (Exp 2-2)}
\label{sec:exp22}

이 절의 목적은 supervised 차원 축소(LDA)가 PCA를 능가할 수 있는지를 검증하는 것이다.

GPT-2 Medium 24층에서 두 가지 비트 설정(6/2-bit, 8/2-bit)으로 PCA와 LDA의 Key MSE를 비교하였다.
Table~\ref{tab:exp22_main}은 핵심 비교 결과를 제시한다.

\begin{table}[H]
\centering
\caption{PCA vs LDA 비트 할당 비교 (Exp 2-2). 출처: \texttt{exp2\_2\_results.json}, \texttt{comparisons} 필드.}
\label{tab:exp22_main}
\begin{tabular}{lrrrrr}
\toprule
\textbf{Config} & \textbf{PCA K MSE} & \textbf{LDA K MSE} & \textbf{Ratio} & \textbf{$t$-stat} & \textbf{$p$-value} \\
\midrule
6/2-bit & 68.3 & 1779.7 & $26.0\times$ & $-6.63$ & $9.14 \times 10^{-7}$ \\
8/2-bit & 59.8 & 1790.1 & $29.9\times$ & $-6.68$ & $8.11 \times 10^{-7}$ \\
\bottomrule
\end{tabular}
\end{table}

Table~\ref{tab:exp22_main}이 보여주듯이, LDA 축은 K MSE에서 PCA 대비 26--30배 열위이다.
이 차이는 통계적으로 고도로 유의하다($p < 10^{-6}$).
V MSE는 두 방법 간 1\% 이내로 유사하여(PCA 6/2: 199.8, LDA 6/2: 198.7; PCA 8/2: 201.5, LDA 8/2: 200.3), K MSE가 차이의 핵심 동력이다.
KL divergence에서는 8/2-bit 설정에서 PCA(1.83)가 LDA(11.87)를 6.5배 능가하나, 6/2-bit에서는 양자 모두 약 12.3--12.4로 유사하다.

이 결과는 LDA가 클래스 분리를 최적화하면서 Key 분포의 분산 구조를 파괴하기 때문으로 해석된다.
양자화 오차 최소화에는 최대 분산 방향의 보존이 핵심이며, 이는 PCA가 정확히 수행하는 작업이다.

\subsection{AFOD vs PCA: 통계적으로 구분 불가 (Exp 2-4)}
\label{sec:exp24}

이 절의 목적은 AFOD(Attention-aware Frequency-Ordered Decomposition) 초기화가 PCA를 개선하는지를 검증하는 것이다.

Table~\ref{tab:exp24_main}은 네 가지 방법(PCA, AFOD, Random, Identity)의 MSE를 비트 폭별로 비교한다.

\begin{table}[H]
\centering
\caption{AFOD vs PCA vs Random vs Identity MSE 비교 (Exp 2-4). 출처: \texttt{exp2\_4\_results.json}, \texttt{summary} 필드.}
\label{tab:exp24_main}
\begin{tabular}{lrrrr}
\toprule
\textbf{Bits} & \textbf{PCA} & \textbf{AFOD} & \textbf{Random} & \textbf{Identity} \\
\midrule
2 & 399.1 & 402.1 & 2239.1 & 2086.6 \\
3 & 285.2 & 285.0 & 1260.9 & 1151.6 \\
4 & 194.3 & 194.2 & 448.2 & 431.4 \\
\bottomrule
\end{tabular}
\end{table}

Table~\ref{tab:exp24_main}에서 PCA와 AFOD의 MSE 차이는 모든 비트 폭에서 1\% 미만이다.
통계 검정 결과, 2-bit($t = 0.70$, $p = 0.49$), 3-bit($t = -0.84$, $p = 0.41$), 4-bit($t = -0.49$, $p = 0.63$)에서 모두 유의한 차이가 없다.
반면, PCA는 Random 대비 모든 비트 폭에서 유의하게 우수하다($p < 1.1 \times 10^{-4}$).
AFOD가 PCA의 개선을 제공하지 못한다는 사실은, PCA 축이 이미 양자화 오차의 하한에 근접해 있음을 시사한다.

\subsection{Lie Group 수치 최적화 vs PCA: PCA 4.3배 우위 (Exp 3-1)}
\label{sec:exp31}

이 절의 목적은 SO($d$) 연속 회전 최적화를 통해 PCA보다 나은 양자화 축을 찾을 수 있는지를 검증하는 것이다.

64개 Givens 회전 평면을 200회 반복 최적화하여, PCA, Identity, Random, Optimized 네 방법의 MSE를 GPT-2 Medium 8개 층에서 비교하였다.
Table~\ref{tab:exp31_rank}는 전체 순위를 제시한다.

\begin{table}[H]
\centering
\caption{양자화 축 최적화 방법 순위 (Exp 3-1). 출처: \texttt{exp3\_1\_results.json}, \texttt{ranking} 필드.}
\label{tab:exp31_rank}
\begin{tabular}{clr}
\toprule
\textbf{Rank} & \textbf{Method} & \textbf{Mean MSE} \\
\midrule
1 & PCA & 220.5 \\
2 & Optimized & 938.8 \\
3 & Identity & 961.0 \\
4 & Random & 1047.2 \\
\bottomrule
\end{tabular}
\end{table}

Table~\ref{tab:exp31_rank}에서 PCA의 평균 MSE(220.5)는 Optimized(938.8)의 4.3배 우수하다.
Optimized vs PCA 비교는 통계적으로 유의하게 PCA가 우월하다($t = -3.45$, $p = 0.011$).
Optimized는 Identity 대비 개선 폭이 층별 0.5--4.9\%에 불과하며, Identity 대비 통계적 유의성도 확보하지 못하였다($t = 2.07$, $p = 0.078$).
PCA는 Random 대비에서도 유의하게 우수하다($t = 3.68$, $p = 0.008$).

이 결과의 의미는 결정적이다.
PCA 기저는 단순히 좋은 선택이 아니라, 수치 최적화가 1시간 이상 탐색해도 도달하지 못하는 수준의 성능을 제공한다.
64개 Givens 평면이라는 제한된 탐색 공간이 원인일 수 있으나, 실용적 관점에서 PCA의 지배적 지위는 확고하다.

\subsection{SSM Distillation: 큰 State Dimension 필요 (Exp 3-2)}
\label{sec:exp32}

이 절의 목적은 attention 패턴을 SSM(State Space Model)으로 증류할 때 필요한 state dimension을 정량화하는 것이다.

Table~\ref{tab:exp32_main}은 state dimension($N$)별 평균 NMSE를 제시한다.

\begin{table}[H]
\centering
\caption{SSM Distillation: State Dimension별 NMSE (Exp 3-2). 출처: \texttt{exp3\_2\_results.json}.}
\label{tab:exp32_main}
\begin{tabular}{rrr}
\toprule
\textbf{$N$} & \textbf{Mean NMSE} & \textbf{Layers $< 0.1$} \\
\midrule
16 & 0.489 & 0/6 \\
32 & 0.454 & 0/6 \\
64 & 0.403 & 0/6 \\
128 & 0.252 & 0/6 \\
256 & 0.157 & 4/6 \\
\bottomrule
\end{tabular}
\end{table}

Table~\ref{tab:exp32_main}에서 NMSE $< 0.1$을 달성하는 층은 $N = 256$에서만 4/6개 존재한다.
$N = 128$에서도 0/6이며, $N = 64$ 이하에서는 평균 NMSE가 0.4 이상으로 실용 가치가 없다.
Head dimension이 64인 GPT-2 Medium에서 $N = 256$은 head dimension의 4배에 해당하여, SSM이 attention 패턴을 압축하기보다 오히려 팽창시키는 결과가 된다.
이 발견은 PCA의 위상을 간접적으로 강화한다. PCA는 추가 파라미터 없이 O($d^2$) 연산만으로 최적 축을 제공하는 반면, SSM은 $N = 256$ 규모의 hidden state를 요구하면서도 완전한 재구성에 실패한다.

\subsection{종합: PCA 최적성의 수렴 증거}

네 실험의 증거를 종합하면, PCA는 양자화 축 선택에서 사실상(de facto) 최적이다.
LDA, AFOD, Lie Group 수치 최적화, SSM 증류 중 어떤 대안도 PCA를 능가하지 못하였다.
이 결론의 실용적 함의는 명확하다.
FOKVQ 양자화 파이프라인에서 축 선택 단계는 PCA 계산 한 번으로 충분하며, 추가적인 학습 기반(AFOD, SSM) 또는 최적화 기반(Givens) 접근은 불필요한 복잡성만 추가한다.

%% ==========================================================================
%% Section 3: Central Finding 2 --- K-space Facet 구조의 실재
%% ==========================================================================
\section{Central Finding 2: K-space Facet 구조의 실재}

이 절은 Key 벡터 공간에 코퍼스 의존적 facet 구조가 존재한다는 발견을 제시하고, 이 구조가 양자화 설계에 미치는 함의를 논한다.

\subsection{Facet K-space 분리도 (Exp 2-1)}
\label{sec:exp21}

이 절의 목적은 서로 다른 코퍼스(technical, code, conversational)가 K-space에서 분리 가능한 하위 공간을 형성하는지를 정량적으로 검증하는 것이다.

GPT-2 Medium 24층에서 코퍼스별 상위 8개 주성분 간 cosine similarity를 측정하였다.
Table~\ref{tab:exp21_main}은 within-corpus와 cross-corpus similarity를 비교한다.

\begin{table}[H]
\centering
\caption{K-space Facet 분리도 (Exp 2-1). 출처: \texttt{exp2\_1\_results.json}.}
\label{tab:exp21_main}
\begin{tabular}{lrl}
\toprule
\textbf{Metric} & \textbf{Value} & \textbf{Source Field} \\
\midrule
Within-corpus similarity (mean) & 0.819 & \texttt{overall\_within\_mean} \\
Cross-corpus similarity (mean) & 0.411 & \texttt{overall\_cross\_mean} \\
Separability gap & 0.408 & \texttt{separability\_gap} \\
$t$-statistic & 30.50 & \texttt{ttest\_stat} \\
$p$-value & $1.64 \times 10^{-64}$ & \texttt{ttest\_p\_value} \\
Mann-Whitney $U$ & 5177 & \texttt{mannwhitney\_stat} \\
Mann-Whitney $p$ & $2.68 \times 10^{-25}$ & \texttt{mannwhitney\_p\_value} \\
\bottomrule
\end{tabular}
\end{table}

Table~\ref{tab:exp21_main}에서 separability gap 0.408은 within-corpus similarity(0.819)의 약 절반에 해당하며, 두 검정 모두 극도로 유의하다.
이 분리도는 하위 층보다 중간--상위 층에서 더 강하게 나타난다.
예를 들어, code vs conversational의 층별 유사도는 L0에서 0.515에서 L7에서 0.280까지 하락한 후, L23에서 0.362로 부분 회복된다.

이 발견의 의미는 이중적이다.
한편으로, 코퍼스 의존적 facet 구조가 실재한다는 것은 모델이 입력 도메인에 따라 Key 공간의 서로 다른 방향을 활성화함을 의미한다.
다른 한편으로, Section~\ref{sec:exp22}--\ref{sec:exp31}에서 보았듯이, 이 구조를 활용한 축 선택(LDA, AFOD)이 PCA를 능가하지 못하므로, facet 구조의 존재가 곧바로 양자화 개선으로 이어지지는 않는다.

\subsection{GQA에 의한 열적 구조 교란 (Exp 1-1)}
\label{sec:exp11_facet}

이 절의 목적은 GQA(Grouped Query Attention)가 K-space의 facet 구조에 어떤 영향을 미치는지를 유효 온도 분석을 통해 간접적으로 확인하는 것이다.

Qwen2.5-3B(GQA, KV heads 2, query heads 16)에서 층별 $\teff$ 평균과 head variance를 분석하였다.
GPT-2 Medium(MHA, heads 16)에서 관찰된 progressive cooling($\rho = -0.657$, $p = 4.82 \times 10^{-4}$)이 Qwen에서는 완전히 사라진다($\rho = +0.108$, $p = 0.530$).

GQA의 구조적 효과는 층 5(Layer 5)에서 극단적으로 나타난다.
Layer 5의 16개 head $\teff$ 값은 최솟값 0.00015(H10)에서 최댓값 0.728(H6)까지 분포한다.
KV Group 0(H0--H7)의 평균 $\teff$는 0.548이고(H5 포함), KV Group 1(H8--H15)의 평균 $\teff$는 0.042이다.
두 그룹 간 gap은 0.506으로, 동일 층 내에서도 KV 그룹에 따라 attention 분포가 극단적으로 분화된다.
Head variance 최댓값은 Layer 5에서 0.0857이며, 이는 Layer 0의 0.0054 대비 15.9배에 달한다.
전체 층의 head variance 중앙값은 0.037이다.

이 관찰은 facet 구조의 해석에 중요한 제약을 부과한다.
GQA 모델에서는 소수의 KV head가 다수의 query head를 서비스하므로, Key 공간의 구조가 MHA와 질적으로 다를 수 있다.
Facet K-space 분리도 실험(2-1)이 MHA 모델(GPT-2)에서 수행되었으므로, GQA 모델에서의 재검증이 필요하다.

\subsection{종합: Facet 구조의 존재와 한계}

K-space facet 구조는 실재하며 통계적으로 강건하다.
그러나 이 구조를 양자화 축 선택에 활용하려는 시도(LDA, AFOD)는 PCA를 능가하지 못하였다.
실용적 결론은 다음과 같다.
Facet 구조는 모델 해석(interpretability)에는 가치가 있으나, 양자화 파이프라인 설계에서는 정적 PCA로 충분하다.

%% ==========================================================================
%% Section 4: Central Finding 3 --- Boltzmann 프레임워크의 가치와 한계
%% ==========================================================================
\section{Central Finding 3: Boltzmann 프레임워크의 기술적 가치와 한계}

이 절은 Boltzmann 유효 온도 프레임워크가 attention 패턴을 기술하는 데 어느 정도까지 유효한지를 세 실험(1-1, 1-2, 1-3)의 결과로 평가한다.

\subsection{Progressive Cooling: MHA에서 확인, GQA에서 파괴 (Exp 1-1)}
\label{sec:exp11}

이 절의 목적은 attention 분포의 유효 온도 $\teff$가 층이 깊어짐에 따라 단조 감소하는지(progressive cooling 가설)를 검증하는 것이다.

Table~\ref{tab:exp11_main}은 두 모델의 핵심 결과를 비교한다.

\begin{table}[H]
\centering
\caption{Progressive Cooling 검증 (Exp 1-1). 출처: \texttt{exp1\_1\_results.json} (Qwen), \texttt{exp1\_1\_gpt2\_results.json} (GPT-2).}
\label{tab:exp11_main}
\begin{tabular}{llrrl}
\toprule
\textbf{Model} & \textbf{Arch} & \textbf{Spearman $\rho$} & \textbf{$p$-value} & \textbf{Verdict} \\
\midrule
GPT-2 Medium & MHA (16 heads) & $-0.657$ & $4.82 \times 10^{-4}$ & \textcolor{partial}{Weak Success} \\
Qwen2.5-3B & GQA (2 KV heads) & $+0.108$ & $0.530$ & \textcolor{failure}{Failure} \\
\bottomrule
\end{tabular}
\end{table}

Table~\ref{tab:exp11_main}에서 GPT-2의 $\rho = -0.657$은 중간 강도의 부적 상관으로, 깊은 층일수록 $\teff$가 감소하는 경향을 확인한다.
GPT-2의 전역 최솟값은 L19($\teff = 0.221$)이며, 최종 층(L23)에서 0.520으로 reheat 현상이 관찰된다.

Qwen에서는 $\rho = +0.108$로, progressive cooling이 완전히 부재한다.
이 차이의 주된 원인은 GQA 구조에 있다.
2개의 KV head가 16개의 query head를 서비스하는 구조에서, 층 내 head 간 $\teff$ 분산이 극대화되어(최대 head variance 0.0857, Layer 5) 층별 평균 $\teff$의 monotonic 감소를 파괴한다.

이 결과는 Boltzmann $\teff$ 프레임워크의 적용 범위를 제한한다.
Progressive cooling은 MHA에 국한된 현상이며, 현대 LLM의 주류인 GQA 아키텍처에서는 성립하지 않는다.

\subsection{비열 $\Cq$와 키 노름: 39.8\% 불일치 (Exp 1-2)}
\label{sec:exp12}

이 절의 목적은 Boltzmann 비열 $\Cq$가 키 노름 $\knorm$와 어느 정도 상관하며, 양자화 비트 할당에 $\Cq$가 기존 variance 기반 방법보다 나은지를 검증하는 것이다.

GPT-2 Medium에서 $\Cq$와 $\knorm$의 층별 상관 분석과 PPL 기반 양자화 비교를 수행하였다.
$\Cq$와 $\knorm$ 간 평균 불일치율(disagreement)은 39.8\%이며, 평균 상관(Spearman $\rho$)은 0.275이다.
층별 상관은 초기 층에서 강하고 후기 층에서 급격히 약화된다.
Early group(L0--L5) 평균 $\rho = 0.526$인 반면, Late group(L18--L23) 평균 $\rho = 0.136$이다.
Mid-Late group(L12--L17) 평균 $\rho$는 0.212이며, Middle group(L6--L11) 평균 $\rho$는 0.226이다.

Table~\ref{tab:exp12_ppl}은 비트 폭별 PPL 비교를 제시한다.

\begin{table}[H]
\centering
\caption{비열 기반 양자화 PPL 비교 (Exp 1-2). 출처: \texttt{exp1\_2\_results.json}, \texttt{ppl\_results} 필드.}
\label{tab:exp12_ppl}
\begin{tabular}{lrrr}
\toprule
\textbf{Config} & \textbf{Uniform} & \textbf{Variance} & \textbf{Specific Heat} \\
\midrule
Baseline & \multicolumn{3}{c}{139.85} \\
2-bit & 419.41 & $4.85 \times 10^{8}$ & $4.85 \times 10^{8}$ \\
3-bit & 130.38 & 127.62 & 264.48 \\
4-bit & 110.19 & 122.85 & 125.56 \\
\bottomrule
\end{tabular}
\end{table}

Table~\ref{tab:exp12_ppl}에서 비열 기반 할당은 어떤 비트 폭에서도 uniform이나 variance 기반 방법을 능가하지 못한다.
3-bit에서 비열(264.48)은 uniform(130.38)의 2배, variance(127.62)의 2.1배에 달하며, 4-bit에서도 비열(125.56)이 uniform(110.19)보다 높다.
Variance 기반 방법이 2승(heat wins = 0, var wins = 2)으로, $\Cq$ 기반 접근은 실용적 가치가 없다.

이 실험의 함의는 Boltzmann 프레임워크의 한계를 명확히 한다.
$\Cq$는 $\knorm$와 독립적인 물리량으로서 기술적 흥미가 있으나, 양자화 비트 할당이라는 처방적 과제에서는 단순한 분산 기반 방법보다 열등하다.

\subsection{FFN 주파수 선택성: 기각 (Exp 1-3)}
\label{sec:exp13}

이 절의 목적은 FFN 블록이 attention의 주파수 선택성(entropy)을 강화하는지를 검증하는 것이다.

GPT-2 Medium에서 FFN 유무에 따른 attention entropy를 24개 층에서 비교하였다.
FFN 존재 시 평균 entropy는 3.376, FFN 제거 시 3.255이며, 평균 차이 $\Delta H = -0.122$이다.
FFN은 entropy를 감소시키는 것이 아니라 증가시킨다($t = -11.79$, $p = 3.15 \times 10^{-11}$; Wilcoxon $W = 0.0$, $p = 2.70 \times 10^{-5}$).
24개 층 중 FFN이 entropy를 감소시키는 층은 0개이다.

이 결과는 ``FFN이 attention의 주파수 선택성을 강화한다''는 가설을 명확히 기각한다.
FFN은 오히려 entropy를 증가시켜 attention 분포를 더 균일하게 만드는 효과를 가진다.
이는 FFN 블록의 역할이 주파수 필터링이 아닌, 표현 공간의 비선형 변환에 있음을 시사한다.

\subsection{종합: 기술적 가치, 처방적 가치 부재}

Boltzmann 프레임워크의 세 실험을 종합하면 다음과 같이 판정된다.
$\teff$는 MHA에서 progressive cooling이라는 실재 현상을 포착하나, GQA에서 파괴된다.
$\Cq$는 독립적 물리량이나 양자화 개선에 실패한다.
FFN 주파수 선택성은 기각된다.
따라서 Boltzmann 프레임워크는 기술적(descriptive) 분석 도구로서 가치를 가지되, 양자화 설계를 위한 처방적(prescriptive) 가치는 입증되지 않았다.

%% ==========================================================================
%% Section 5: Negative/Inconclusive 결과 정리
%% ==========================================================================
\section{Negative 및 Inconclusive 결과 정리}

이 절은 가설이 기각되었거나 결론을 내릴 수 없었던 실험들을 정리하고, 재실험 시 개선 방향을 제시한다.

\subsection{Exp 2-3: Query-Dependent vs Static 할당 --- Inconclusive}
\label{sec:exp23}

이 절의 목적은 query-dependent 동적 비트 할당이 정적 PCA 할당을 능가하는지를 검증하는 것이었다.

이 실험은 \textbf{Inconclusive}로 판정한다.
두 독립 구현(A, B)이 상충하는 결과를 산출하였기 때문이다.
구현 A는 \texttt{query\_dependent\_allocation\_differs = false}와 NaN으로 가득한 rank correlation을 보고하여 ``정적 할당이 충분하다''는 결론을 지지하였다.
구현 B는 \texttt{query\_dependent\_allocation\_differs = true}와 유효한 rank correlation($\rho \approx 0.70$)을 보고하여 반대 결론을 지지하였다.
현재 JSON 파일(\texttt{exp2\_3\_results.json})에는 구현 B의 결과가 저장되어 있으나, 구현 A의 데이터가 먼저 존재하였고 구현 B가 이를 덮어쓴 경과가 있다.

어느 구현이 정확한지를 현 데이터만으로 판별할 수 없으므로, 양쪽의 구체적 수치를 인용하지 않는다.
이 실험은 다음 두 조건이 충족될 때 재수행되어야 한다.
첫째, per-dimension importance와 per-head importance의 구현이 분리되어 각각 독립적으로 검증되어야 한다.
둘째, degenerate allocation(모든 차원이 동일 bucket에 배치되는 경우)을 사전에 탐지하는 sanity check가 추가되어야 한다.

\subsection{Exp 3-3: SSM 초기화 --- Spectral 가설 기각}
\label{sec:exp33}

이 절의 목적은 spectral 초기화가 SSM의 attention 패턴 학습에서 최적인지를 검증하는 것이다.

GPT-2 Medium Layer 12에서 5가지 초기화 방법(HiPPO, Random, RoPE, YaRN, Spectral)을 비교하였다.
Table~\ref{tab:exp33_main}은 최종 NMSE 순위를 제시한다.

\begin{table}[H]
\centering
\caption{SSM 초기화 방법 NMSE 비교 (Exp 3-3). 출처: \texttt{exp3\_3\_results.json}, \texttt{ranking} 필드.}
\label{tab:exp33_main}
\begin{tabular}{clrr}
\toprule
\textbf{Rank} & \textbf{Method} & \textbf{NMSE} & \textbf{Best Loss} \\
\midrule
1 & RoPE & 0.475 & 0.410 \\
2 & Spectral & 0.496 & 0.434 \\
3 & Random & 0.500 & 0.428 \\
4 & YaRN & 0.513 & 0.438 \\
5 & HiPPO & 0.621 & 0.527 \\
\bottomrule
\end{tabular}
\end{table}

Table~\ref{tab:exp33_main}에서 RoPE 초기화(NMSE = 0.475)가 최우수이며, Spectral(0.496)은 2위이나 Random(0.500)과 거의 구분되지 않는다.
HiPPO 초기화가 최하위(0.621)인 것은, 이론적으로 장기 기억에 최적화된 HiPPO 행렬이 attention의 local 패턴과 부합하지 않기 때문으로 해석된다.
어떤 방법도 수렴(convergence)을 달성하지 못하였으며, 유일한 예외는 YaRN(\texttt{converged = true})이나 NMSE 자체는 4위이다.

이 결과는 두 가지를 시사한다.
첫째, Spectral 초기화가 최적이라는 가설은 기각된다. RoPE가 더 나은 초기 조건을 제공한다.
둘째, 모든 방법의 NMSE가 0.47--0.62 범위에서 정체되어 있어, 300 step의 학습으로는 attention 패턴을 충분히 재구성할 수 없다.
재실험 시에는 학습 step 수를 최소 1000 이상으로 확대하고, state dimension 64 이상(Exp 3-2 결과 참조)을 사용하여야 한다.

\subsection{Negative 결과의 과학적 기여}

Negative 결과는 탐색 공간을 축소한다는 점에서 과학적 가치를 가진다.
FFN 주파수 선택성 기각(1-3)은 향후 양자화 연구에서 FFN 통합 접근을 배제할 근거를 제공하고, Spectral 초기화 기각(3-3)은 SSM 기반 attention 압축 연구에서 RoPE 초기화를 기본값으로 채택할 근거를 제공한다.
Exp 2-3의 inconclusive 결과는 구현 검증의 중요성을 환기시킨다.

%% ==========================================================================
%% Section 6: 종합 판정과 Next Steps
%% ==========================================================================
\section{종합 판정과 Next Steps}

이 절은 10개 실험의 전체 결론을 제시하고, Phase 4 확장 실험의 방향을 제안한다.

\subsection{종합 판정}

10개 실험의 결론은 세 문장으로 요약된다.
첫째, 양자화 축 선택에서 PCA는 모든 대안(LDA, AFOD, Lie Group, SSM)을 압도하며, 추가적 최적화 시도는 불필요하다.
둘째, K-space facet 구조는 실재하나(gap = 0.408), 양자화 설계에서의 직접적 활용 가치는 아직 입증되지 않았다.
셋째, Boltzmann $\teff$ 프레임워크는 MHA에서 기술적 유효성을 보이나, GQA에서 파괴되고 처방적 가치가 없어 현재 형태로는 양자화 설계 도구로 부적합하다.

\subsection{Phase 4 로드맵}

Table~\ref{tab:phase4}는 제안하는 Phase 4 실험을 정리한다.

\begin{table}[H]
\centering
\caption{Phase 4 실험 제안.}
\label{tab:phase4}
\begin{tabularx}{\textwidth}{lX}
\toprule
\textbf{실험} & \textbf{내용} \\
\midrule
4-1 & Llama-3-8B(GQA)에서 Exp 2-1 재현: GQA 모델에서의 facet 분리도 검증 \\
4-2 & Exp 2-3 재실험: per-dimension importance 구현 검증 후 단일 정본 결과 산출 \\
4-3 & PCA 축 + calibration set 크기 민감도 분석: 최소 calibration 크기 결정 \\
4-4 & Mixed-precision PCA: 층별 비트 폭 자동 결정(NMSE threshold 기반) \\
\bottomrule
\end{tabularx}
\end{table}

\subsection{논문 구성 제안}

본 실험 시리즈를 학술 논문으로 발전시킬 경우, 다음 구성을 제안한다.
논문의 핵심 주장(central claim)은 ``PCA가 KV cache 양자화의 최적 축 선택이다''로 설정한다.
Introduction에서 Boltzmann 프레임워크를 동기 부여(motivation)로 제시하되, 본론에서는 PCA 최적성의 증거(Exp 2-2, 2-4, 3-1)를 중심 축으로 구성한다.
Facet 구조(Exp 2-1)는 PCA가 왜 잘 작동하는지에 대한 해석(interpretation)으로 활용한다.
Negative 결과(Exp 1-2, 1-3, 3-2, 3-3)는 Ablation Study 섹션에 배치하여, PCA 대안이 실패하는 조건을 체계적으로 보여준다.
Exp 2-3의 inconclusive 결과는 Limitations 섹션에서 향후 과제로 명시한다.

%% ==========================================================================
%% Appendix: Data Source Reference
%% ==========================================================================
\appendix
\section{데이터 출처 참조표}

본 보고서의 모든 수치는 다음 JSON 파일에서 직접 추출하였다.

\begin{table}[H]
\centering
\caption{실험별 데이터 파일 및 주요 필드.}
\label{tab:datasource}
\small
\begin{tabularx}{\textwidth}{lXl}
\toprule
\textbf{Exp} & \textbf{File} & \textbf{Key Fields} \\
\midrule
1-1a & \texttt{exp1\_1\_results.json} & \texttt{spearman\_rho}, \texttt{head\_var\_per\_layer} \\
1-1b & \texttt{exp1\_1\_gpt2\_results.json} & \texttt{spearman\_rho}, \texttt{layer\_teff\_mean} \\
1-2 & \texttt{exp1\_2\_results.json} & \texttt{avg\_disagreement}, \texttt{ppl\_results} \\
1-3 & \texttt{exp1\_3\_results.json} & \texttt{mean\_diff}, \texttt{paired\_t\_stat} \\
2-1 & \texttt{exp2\_1\_results.json} & \texttt{separability\_gap}, \texttt{ttest\_stat} \\
2-2 & \texttt{exp2\_2\_results.json} & \texttt{comparisons}, \texttt{per\_config\_results} \\
2-3 & \texttt{exp2\_3\_results.json} & (Inconclusive; 두 구현 상충) \\
2-4 & \texttt{exp2\_4\_results.json} & \texttt{summary}, \texttt{statistical\_tests} \\
3-1 & \texttt{exp3\_1\_results.json} & \texttt{ranking}, \texttt{statistical\_tests} \\
3-2 & \texttt{exp3\_2\_results.json} & \texttt{N*\_mean\_nmse} \\
3-3 & \texttt{exp3\_3\_results.json} & \texttt{ranking}, \texttt{methods} \\
\bottomrule
\end{tabularx}
\end{table}

\section{Critic Corrections 적용 내역}

이 보고서는 Critic Rounds 1--5에서 지적된 모든 오류를 반영하였다.

\begin{table}[H]
\centering
\caption{Critic Corrections 적용 목록. 출처: \texttt{critic\_corrections.md}.}
\label{tab:critic}
\small
\begin{tabularx}{\textwidth}{llX}
\toprule
\textbf{ID} & \textbf{원본 오류} & \textbf{수정 내용} \\
\midrule
R1 & HeadVar middle = 0.036 & 0.037로 수정 (\S\ref{sec:exp11_facet}) \\
R2 & Exp 2-2 KL ``PCA=4.75, LDA=13.28'' & PCA=1.83, LDA=11.87로 수정 (\S\ref{sec:exp22}) \\
F1 & ``KV0 평균 = 0.548 (H5 제외 시)'' & ``H5 포함'' 으로 수정 (\S\ref{sec:exp11_facet}) \\
F2 & Mid-Late $\rho$ = 0.183 & 0.212로 수정 (\S\ref{sec:exp12}) \\
F3 & Middle $\rho$ = 0.219 & 0.226으로 수정 (\S\ref{sec:exp12}) \\
F4 & FOKVQ .md 인용 & .docx 출처 미검증; 본 보고서에서 인용 제외 \\
F5 & Exp 2-3 old data 기반 서술 & Inconclusive로 재판정, 양쪽 수치 비인용 (\S\ref{sec:exp23}) \\
\bottomrule
\end{tabularx}
\end{table}

\end{document}
